% Irodov Part V §5.4 Polarization of Light, original problem 157.
% Original uploaded PDF p.222 (printed p.213), answers p.332.
\polarizationproblem{157}{213}
强度为 $I_0$ 的单色\emph{自然光}平面波垂直照到一块屏幕，
屏幕由两片在直边处相接的理想偏振片半平面组成：
一片透振方向平行于分界线，
另一片透振方向垂直于分界线。
问：(a) 屏后形成什么衍射图样；
(b) 位于经过两偏振片边界且与屏垂直的平面内各观察点的光强。
\exsolution
自然光可分为统计独立、强度各为 $I_0/2$ 的两正交线偏振分量。
对于透振方向沿边界的分量，只有第一半平面透过；
对于另一正交分量，只有第二半平面透过。
两半平面的输出电场始终正交，因而它们的\emph{强度可相加，电场不能产生互相干涉项}。
每一偏振分量在边界附近形成普通菲涅耳半平面衍射强度分布，
两列分布互为边界镜像：
\[
\boxed{I(x)=\frac{I_0}{2}
\left[\mathcal H(x)+\mathcal H(-x)\right]},
\]
其中 $\mathcal H(x)$ 表示以相应透射半平面入射强度1归一化的菲涅耳边缘衍射强度函数。
对位于分界线法平面的边界方向，
两分量各受半平面遮挡，场振幅为全开放场的一半，
故其各自强度为
$(I_0/2)(1/2)^2=I_0/8$。
因此
\[
\boxed{I_{\rm boundary}=
\frac{I_0}{8}+\frac{I_0}{8}
=\frac{I_0}{4}}.
\]
\excheck 两半偏振片透振方向互相正交，
因此不会出现两半面之间的普通相干明暗干涉条纹；
边缘附近仍分别存在菲涅耳\emph{衍射}，这与“没有衍射”完全不同。
